\paragraph{R2-W2, R2-D2 (Missing critical evidence; empty-set fallback).}

\begin{itemize}[leftmargin=1.2em, itemsep=1pt, topsep=2pt]
\item \textit{W2: ``what about missing critical evidence for the counterfactual? The core premise of a data-centric approach requires observed instances that support alternatives.''}
\item \textit{D2: ``The data-central approach requires the counterfactual evidence to be present in $\mathcal{D}$. Some discussions around this (e.g., empty-set fallback) would be interesting.''}
\end{itemize}
\smallskip

\noindent\textbf{Response.}
By Eq.~(1), an instance is admissible \emph{``if and only if $\omega(\mathbf{x}_q,\mathbf{x})<\infty$''}; when $\omega(q)=\varnothing$, DcCA returns an empty set that \emph{certifies} no feasible, outcome-changing real instance exists---a verifiable statement generative methods cannot make~\cite{wachter2017counterfactual,mothilal2020explaining}. Our validity study quantifies this: in Table~\ref{tab:validity}, a feasible immutability-matched counterfactual exists for $83.7\%$ of queries, and the remaining $16.3\%$ fall into exactly this empty-set regime. We add a graded, \emph{data-centric} fallback: the certified empty set, partial sets with provenance, and relaxation of soft constraints reporting the minimal relaxation that restores feasibility. Data-centric and model-centric analysis answer different questions rather than serving as fallbacks for one another, so DcCA does not depend on a generative method as a fallback.

\noindent\emph{Revision to Section~2 (new paragraph ``Evidence availability and query semantics''):}
\begin{quote}
\textcolor{blue}{When $\omega(q)=\varnothing$, no feasible, outcome-changing instance exists in $\mathcal{D}$, and DcCA returns the empty set, certifying that no evidence-supported counterfactual exists under $\omega$. For $0<|\omega(q)|<k$ it returns all feasible counterfactuals with their provenance. Data-centric and model-centric analysis answer different questions rather than serving as fallbacks for one another: data-centric retrieval guarantees realism, feasibility, and verifiability but is bounded by the coverage of $\mathcal{D}$, whereas generative methods always return an answer at the cost of feasibility and verifiability. Coverage is measurable and grows with $\mathcal{D}$.}
\end{quote}
